% !TeX root = ../samplepaper.tex
\section{Conclusion}
本文从local search的候选起点出发设计FACO外层控制，将参考路径切换、扰动起始区域和强度联合表述为GP可学习的决策。在固定构造与local-search组件下，TSP500上的两种表示均学到优于MNE=8配置的规则，但MNE=16静态对照取得接近质量，TSP1000上的区域静态策略进一步超过两种GP。规则分析与采样决策显示对MNE=16的偏好，反馈主要影响区域和重启动作。质量改善伴随着更多LS工作；multi-tree也未以更高训练成本换来明确质量收益。

这些结果支持围绕local search考察外层控制，同时将尚待验证的问题集中到状态相关决策的增量价值。后续应在另行划定的预实验实例集固定MNE=16，对反馈、区域选择及参考路径切换进行完整求解消融，联合测量终点质量与LS工作量，并检查停滞特征、重启频率及训练与部署预算的差异。现有scalar fitness由多个决策共同承担credit assignment，动作敏感性还不足以分离各项控制的质量贡献。结论来自三个GP seed、TSP500训练与TSP1000迁移；FACO面向大规模问题的设计背景不等同于本方法已完成大规模泛化验证。
